\chapter{Présentation générale et cadrage du projet}
\label{chap:generalites}

\section*{Introduction}
La transformation numérique des activités commerciales impose aux entreprises de mieux structurer le suivi de leurs clients, de leurs créances et des actions de recouvrement. Dans ce contexte, la qualité des outils internes influence directement la réactivité des équipes, la traçabilité des opérations et la fiabilité des données exploitées pour la prise de décision.

Ce chapitre présente le cadre général du projet \textbf{Collectra}, réalisé dans le cadre d'un stage de Projet de Fin d'Études. Il expose le contexte, la problématique, les objectifs, le périmètre fonctionnel et les choix techniques globaux qui orientent le reste du rapport.

\section{Contexte général}
Le recouvrement et la gestion des dettes clients reposent souvent sur des processus hétérogènes (tableurs, échanges manuels, outils non unifiés), ce qui entraîne :
\begin{itemize}
  \item une dispersion des informations client ;
  \item un suivi difficile de l'historique des actions ;
  \item des risques d'erreurs dans la coordination entre collaborateurs ;
  \item une faible visibilité sur l'avancement réel des dossiers.
\end{itemize}

Le projet \textbf{Collectra} répond à ce besoin par une plateforme web de type SaaS orientée gestion multi-espaces de travail (workspaces), permettant d'organiser les données, de sécuriser l'accès et de standardiser les opérations de suivi.

\section{Présentation du projet Collectra}
\subsection{Nature du projet}
Collectra est une application web full-stack conçue pour centraliser la gestion de :
\begin{itemize}
  \item clients ;
  \item dettes ;
  \item promesses de paiement ;
  \item actions de suivi associées aux dossiers ;
  \item membres d'équipe et rôles au sein d'un espace de travail.
\end{itemize}

L'application adopte une architecture moderne séparant clairement le frontend, l'API métier et la couche d'accès aux données.

\subsection{Positionnement fonctionnel}
Le système vise à offrir :
\begin{itemize}
  \item un tableau de bord privé pour les utilisateurs authentifiés ;
  \item une gestion des workspaces pour isoler les données par contexte métier ;
  \item une administration des membres (invitation, rôle, activation/désactivation) ;
  \item des APIs sécurisées et documentées pour les opérations métier.
\end{itemize}

\section{Problématique}
La problématique principale peut être formulée comme suit :

\begin{quote}
\textit{Comment concevoir une plateforme web sécurisée, multi-tenant et évolutive permettant de gérer efficacement le cycle de suivi des créances, tout en garantissant la traçabilité des actions et la séparation des données entre espaces de travail ?}
\end{quote}

Cette problématique englobe des enjeux à la fois techniques (sécurité, architecture, performance) et organisationnels (collaboration, gouvernance des rôles, continuité du suivi).

\section{Objectifs du projet}
\subsection{Objectif général}
Développer une solution SaaS nommée \textbf{Collectra} permettant la gestion centralisée des clients et des dettes, avec un contrôle d'accès robuste et une organisation multi-workspaces.

\subsection{Objectifs spécifiques}
Les objectifs opérationnels retenus sont :
\begin{itemize}
  \item mettre en place une authentification sécurisée (sessions, jetons, protections API) ;
  \item implémenter la gestion des workspaces et du contexte locataire (tenant) ;
  \item offrir des modules métier pour clients, dettes, promesses et actions ;
  \item permettre la gestion des équipes (invitation, rôles \texttt{MANAGER}/\texttt{AGENT}, statut des membres) ;
  \item assurer la cohérence des données via une base relationnelle et un ORM typé ;
  \item fournir une base logicielle maintenable, testable et extensible.
\end{itemize}

\section{Périmètre du chapitre et du projet}
\subsection{Périmètre fonctionnel couvert}
Le périmètre principal de Collectra inclut :
\begin{itemize}
  \item authentification et gestion des utilisateurs ;
  \item gestion des workspaces et sélection du contexte actif ;
  \item gestion des entités métier (clients, dettes, promesses, actions) ;
  \item gestion des membres internes et des invitations ;
  \item interface web de consultation et d'interaction.
\end{itemize}

\subsection{Hors périmètre à ce stade}
Dans l'état actuel, les éléments suivants ne constituent pas le cœur du livrable :
\begin{itemize}
  \item intégrations avancées à des ERP externes ;
  \item analytique décisionnelle complète (BI) ;
  \item déploiement multi-région à très grande échelle.
\end{itemize}

\section{Aperçu technologique}
Le projet s'appuie sur une architecture monorepo autour des technologies suivantes :
\begin{itemize}
  \item \textbf{Frontend} : Next.js (App Router), React, TypeScript, Tailwind CSS ;
  \item \textbf{Backend} : Hono.js, TypeScript, validation via Zod, documentation OpenAPI ;
  \item \textbf{Données} : PostgreSQL (Supabase) et Prisma ORM ;
  \item \textbf{Auth \& sécurité} : Supabase Auth, JWT, cookies HTTP-only, contrôle d'accès par middleware ;
  \item \textbf{Organisation du code} : Turborepo + pnpm workspaces (applications et packages partagés).
\end{itemize}

Ce socle favorise une bonne séparation des responsabilités, la réutilisabilité des composants et une évolution progressive des fonctionnalités.

\section{Méthodologie de réalisation (vue d'ensemble)}
La réalisation suit une approche incrémentale :
\begin{itemize}
  \item analyse des besoins fonctionnels prioritaires ;
  \item modélisation des données et mise en place des migrations ;
  \item développement des APIs métier sécurisées ;
  \item intégration frontend et gestion d'état ;
  \item validation progressive par tests techniques ciblés.
\end{itemize}

Cette démarche permet de livrer des fonctionnalités exploitables par itérations, tout en maintenant la qualité du code et la cohérence globale du système.

\section*{Conclusion}
Ce premier chapitre a introduit le cadre général du projet Collectra, sa problématique, ses objectifs et son périmètre. Il a également présenté les choix technologiques majeurs et la logique méthodologique adoptée.

Le chapitre suivant détaillera l'analyse des besoins et la conception du système, en mettant l'accent sur les acteurs, les cas d'utilisation et la modélisation des données.
